\section{Discussion}
\label{sec:discussion}

The principal finding is a controlled negative benchmark: the runs completed
validly, but the tested performance hypotheses were rejected. The Q01
fidelity-style quantum kernel was far behind the physics-residual control, and
the Q01b projected feature map did not close the gap. The feasibility-only
kernel did not meet its fixed multi-metric screening criteria. The successor
branch demonstrates an additional failure mode that matters beyond this task: a
candidate can improve over one baseline yet fail against a classical model with
the same staged construction. That distinction is why matched dequantization
controls and thresholds fixed before outcomes belong in the protocol.

The result is scientifically useful because the negative conclusion is
specific. It rules out a family of tested feature maps and successor
constructions under a particular public-data scenario generator, grouped
development procedure, resource envelope, and comparator set. It does not
rule out all QML, all cislunar machine learning, or all future quantum
hardware. Such broad claims would require evidence absent here, including
separate tasks, new preregistration, scalability studies, noise analysis,
hardware measurements, and appropriate classical resource comparisons.

The physics-residual control is a key methodological result. A generic
classical baseline is not sufficient when the prediction target contains
structure that a low-fidelity model can represent directly. C06's performance
shows that, for this benchmark, the available physical residual structure is
more informative than the tested quantum fidelity features. That observation
does not establish a theorem about quantum circuits; it establishes a
task-specific comparison of inductive bias, the kinds of patterns each model is
built to represent, that future proposals must beat fairly.

The retained Q02/Q03 status also matters. Many reports collapse all
non-advancing paths into a single ``failure'' category. Here, the first-rung
tasks ran correctly, but the candidates were ineligible for larger-rung
computation under a rule frozen before results. Keeping
\texttt{not\_reached\_under\_frozen\_eligibility} separate from a task failure
makes the evidence ledger more truthful and prevents a later analyst from
inventing a selection path that never existed.

The paper's societal relevance is therefore not a claim of operational benefit.
It is a reusable, human-centered research pattern for high-consequence AI:
include an eligible simulator-validation scope, freeze data transformations and
groups, use a physics-aware control, compare against matched classical
alternatives, retain null outcomes, and make the claim boundary checkable from
source artifacts. This helps researchers, reviewers, and future decision-tool
designers distinguish a reproducible negative result from evidence for
automation. Its value is in reducing the risk that an exciting but unsupported
model claim is presented as a safe decision-support capability.

The HCI implication is deliberately limited. The present results do not show
that a planner would understand, trust, reject, or use the evidence package
well. They specify what a later human-centered evaluation should not omit:
provenance, model scope, uncertainty, failure visibility, and a non-automated
fallback. Testing those requirements with intended users remains future work.

\subsection{Implications for QML benchmarking}

The QML literature has repeatedly shown that benchmark conclusions depend on
which classical methods are made available to the comparison
\cite{bowles2024,schnabel2024}. The D046 result gives a concrete trajectory
application example. Comparing only with C06 would make the candidate's raw
5.80\% difference look promising. The matched TWO-RBF alternative explains
enough of that difference that the candidate fails the frozen 5\% rule. The
interpretation changes because the question changes from ``is the candidate
better than one baseline?'' to ``is its incremental mechanism still supported
when a classical model is given the same residual construction?''

This contrast does not show that the matched control is the universally best
classical model. It shows why a QML result needs multiple controls chosen to
test the claimed mechanism. A fair comparison is not one that intentionally
makes every model identical; it is one that removes alternative explanations
for the purported QML contribution. In this benchmark, the alternative
explanation was a classical residual representation, and it remained
competitive.

\subsection{Safety objective as a distinct scientific question}

The FQK and CSAFE-RF records demonstrate that predictive metrics must be tied
to both label meaning and intended action. In these executed analyses, the
positive label is feasible. AUROC measures overall ranking, Brier score measures
probability error, and recall measures how many actually feasible options are
retained. Precision and the false-positive rate address the different risk of
admitting an actually infeasible option. FQK remains negative under its frozen
multi-metric rule even though its AUROC is above random, because its Brier score
and feasible-case recall did not match the classical controls. This guards
against promoting a model by whichever metric looks most favorable afterward.

The higher-recall classical result is not a rescue or a safety result. It is a
future design lesson: a screening protocol should define, before fitting, a
minimum feasible-case recall for usefulness, a maximum infeasible-case
false-positive rate for protection, a calibration requirement, and the action
taken when the model abstains. Those quantities cannot be replaced by a single
``safety'' score. Testing such a policy requires a separately authorized
experiment; using it to alter Gate 5 after inspection would invalidate the
original test.
